\chapter{Popularne metody kwantyzacji modeli}

Rozwój dużych modeli językowych (LLM) oraz modeli wizji komputerowej (VLM) doprowadził do sytuacji, w której wdrożenia produkcyjne oraz badania nad architekturami transformerów stały się domeną wyłącznie korporacji dysponujących wydajnymi klastrami obliczeniowymi. Kwantyzacja, polegająca na redukcji precyzji reprezentacji liczb (najczęściej z 32-bitowego formatu zmiennoprzecinkowego do 4- lub 8-bitowych wartości stałoprzecinkowych bądź niestandardowych formatów zmiennoprzecinkowych), stała się kluczowym mechanizmem zwiększającym popularność sztucznej inteligencji. Pozwala ona na drastyczne zmniejszenie zapotrzebowania na przepustowość pamięci oraz pojemność pamięci, przyspieszając wnioskowanie i umożliwiając uruchamianie modeli na sprzęcie konsumenckim.

W niniejszym rozdziale przedstawione zostaną najbardziej rozpowszechnione rozwiązania dziedzinie kwantyzacji. Omówiono zarówno standardy  BitsAndBytes z biblioteki Hugging Face, algorytmy kompresji wag bazujące na analizie błędów (AWQ, GPTQ), nowatorskie formaty plików zoptymalizowane pod heterogeniczne architektury CPU i GPU (GGUF), jak i przełomowe metody mapowania rozkładu wag oparte na transformacjach ortogonalnych (TurboQuant).

\section{BitsAndBytes oraz format NF4}
\label{sec:bnb}

Choć biblioteka BitsAndBytes (BnB) początkowo funkcjonowała jedynie jako biblioteka w ekosystemie PyTorch, dziś nazwa ta utożsamiana jest z de facto standardem w kwantyzacji, głównie dzięki artykułowi \textit{QLoRA} (Dettmers i in., 2023), który spopularyzował format \textbf{NF4} (NormalFloat4). Metoda ta pozwoliła na dostrojenie modeli rzędu 65 miliardów parametrów na pojedynczym akceleratorze graficznym, co wcześniej było niemożliwe.

Standardowe kwantyzacje (np. INT8) zakładają jednostajny rozkład wartości. Oznacza to, że odległość między poszczególnymi poziomami reprezentacji (krok kwantyzacji $\Delta$) jest stały. Badania wykazały jednak, że wytrenowane wagi modeli sieci neuronowych przyjmują rozkład zbliżony do normalnego (o zerowej średniej, skupiony wokół zera z długimi ogonami). W związku z tym przypisanie równych fragmentów pamięci bitowej dla rzadkich wartości na ogonach i gęsto skupionych wartości wokół zera prowadzi do niepotrzebnej straty precyzji (tzw. błąd kwantyzacji).

BitsAndBytes rozwiązuje to za pomocą \textbf{kwantyzacji kwantylowej}. Algorytm wyznacza wartości odczytane z funkcji odwrotnej do skumulowanej funkcji rozkładu normalnego (z ang. Inverse Cumulative Distribution Function - ICDF). W ten sposób uzyskuje się optymalne, \textbf{niestandardowe poziomy wartości}, które gwarantują, że każdy z tzw. "przedziałów" (z ang. bins) posiada dokładnie takie samo prawdopodobieństwo wystąpienia wagi z danego przedziału.

\subsection{Warianty i podtypy w architekturze BitsAndBytes}

Ekosystem BitsAndBytes rozwinął się w stronę wsparcia wielu architektur i potrzeb treningowych, wprowadzając szereg podtypów kwantyzacji:

\begin{itemize}
    \item \textbf{LLM.int8()}: Pionierska metoda umożliwiająca wnioskowanie (inferencję) w formacie 8-bit bez zauważalnego spadku wydajności. Problemem wczesnych modeli była anizotropia aktywacji – podczas wnioskowania pojawiały się tzw. wartości odstające w stanach ukrytych, które mogły mieć amplitudę nawet 100 razy większą niż reszta wymiarów. Standardowa kwantyzacja INT8 niszczyła model. LLM.int8() wprowadziło hybrydową dekompozycję: macierze mnożone są w formacie INT8 z funkcją typu mieszanej precyzji (z ang. "mixed-precision"). Kanały zawierające wartości odstające (zazwyczaj ok. 0.1\% wszystkich wymiarów) izolowane są do mnożenia w pełnej precyzji FP16.
    
    \item \textbf{NF4 (NormalFloat4)}: Główny format używany do kompresji pamięciowej podczas treningu z wykorzystaniem techniki QLoRA. Należy tu podkreślić, że dekwantyzacja z NF4 do formatu BF16/FP16 na potrzeby obliczeń zachodzi w trakcie działania. Oszczędność wynika z faktu, że wagi modelu spoczywające w pamięci VRAM zajmują ułamek oryginalnej wielkości.
    
    \item \textbf{FP4 (Float4)}: Alternatywny format dla NF4. Zamiast niestandardowych przedziałów mieszczących się w rozkładzie normalnym, używa klasycznej reprezentacji zmiennoprzecinkowej z bitem znaku, wykładnikiem i mantysą. Ze względu na bardzo małą liczbę bitów wykładnika, jest on niezwykle podatny na tzw. niedomiar/przepełnienie i jest zazwyczaj gorzej oceniany w testach niż NF4.
    
    \item \textbf{Podwójna Kwantyzacja (Double Quantization - DQ)}: Nawet przy kwantyzacji do 4 bitów potrzebne są stałe normalizujące (skalowanie i bias) zazwyczaj przechowywane w formacie FP32 dla każdej grupy wag. BnB wprowadziło podwójną kwantyzację, która kwantyzuje same stałe (tzw. metadane) z FP32 do FP8, co pozwala zaoszczędzić dodatkowe 0.37 bita na każdą wagę oryginalnego modelu. Co ciekawe, w BnB DQ parametry pomocnicze rozmiaru grupy wynoszą 256, co zachowuje strukturalną integralność macierzy Hesjana procesu uczenia.
\end{itemize}

\section{GPTQ (Accurate Post-Training Quantization)}
\label{sec:gptq}

Algorytm GPTQ, zaproponowany przez Frantara pod koniec 2022 roku, wyznaczył zupełnie nowe standardy dla kwantyzacji po treningu (PTQ - Post-Training Quantization). Jego główną zaletą jest to, że potrafi on zredukować precyzję modeli posiadających setki miliardów parametrów (np. modele takie jak OPT-175B czy BLOOM-176B) w zaledwie kilka godzin obliczeń GPU, przy zachowaniu niemal identycznej dokładności względem modelu oryginalnego.

GPTQ bazuje na koncepcjach algorytmu OBD (Optimal Brain Damage) z początku lat 90., który służył do "przycinania" sieci neuronowych. W skrócie, GPTQ rozwiązuje następujący problem: mając warstwę liniową o wagach $W$, chcemy znaleźć skwantyzowane wagi $\hat{W}$, tak aby błąd rekonstrukcji wyjścia warstwy (błąd aktywacji) był jak najmniejszy dla danego wejścia $X$. Minimalizowana jest funkcja:
\begin{equation}
    \arg\min_{\hat{W}} || WX - \hat{W}X ||^2
\end{equation}

Ideą algorytmu jest kompresja wag wierszami (lub kolumnami w zależności od implementacji) oraz obserwowanie błędu powstającego na wyjściu, by ten sam błąd natychmiast kompensować na wagach sąsiednich, jeszcze nie skwantyzowanych. Obliczanie optymalnych modyfikacji wymaga zróżniczkowania funkcji błędu po wadze względem pozostałych wag z wykorzystaniem Macierzy Hesjana $H_F^{-1}$.

Aby algorytm zadziałał prawidłowo, macierz Hesjana jest statystycznie aproksymowana na małym zbiorze kalibracyjnym – wymagane jest zaledwie kilkaset próbek danych wejściowych. Aby zapobiec załamaniu numerycznemu (inwersja macierzy dla tysięcy wymiarów jest niestabilna), twórcy wdrożyli szybką dekompozycję Cholesky'ego i zjawisko "opóźnionych aktualizacji partii". GPTQ wylicza korekcję w paczkach po 128 kolumn jednocześnie.

\subsection{Warianty i podtypy w ekosystemie GPTQ}

Podczas gdy sam algorytm bazowy GPTQ jest tylko matematycznym modelem PTQ, społeczność open source (zwłaszcza projekt AutoGPTQ) wypracowała cały system podtypów i wariantów dopasowanych do sprzętu.

\begin{itemize}
    \item \textbf{Różne szerokości bitowe}: Najpopularniejszym formatem zaimplementowanym w GPTQ jest kompresja do \textbf{4-bitowych liczb całkowitych (INT4)}. Oferuje on najlepszy stosunek kompresji do spadku "perplexity". Warianty takie jak INT3 i INT8 również są dostępne, jednak INT3 powoduje już zauważalną degradację właściwości logicznych, a INT8 daje zbyt małą kompresję, by uzasadnić narzut dekwantyzacji.
    
    \item \textbf{Rozmiar Grupy}: GPTQ stosuje kwantyzację grupową. Aby równanie kompensacji Hesjana nie okazało się zbyt ogólne (co ma miejsce w kanałowej), a jednocześnie nie pochłaniało nadmiernie pamięci na metadane (jak w przypadku elementowej), grupuje się wagi w bloki. Najpopularniejsze wartości parametru rozmiar grupy to:
    \begin{itemize}
        \item \textbf{-1 (Kanałowa/Kolumnowa)}: Jeden zestaw parametrów (skala, punkt zerowy) dla całego wymiaru wejściowego. Oszczędność pamięci jest największa, ale model jest podatny na anomalię dzielenia przez zero.
        \item \textbf{128}: Standardowa kompresja zbalansowana. Każde 128 wag posiada własny współczynnik skali $s$ i przesunięcie $z$.
        \item \textbf{32}: Bardzo wysoka precyzja dekwantyzacji. Metadane dla grup 32-bitowych potrafią zająć tyle samo miejsca co skwantyzowane wagi, co w rezultacie daje bitrate na poziomie 5.5-6 bitów zamiast założonych 4.
    \end{itemize}
    
    \item \textbf{desc-act (Kolejność aktywacji / statystyki rozkładu GPTQ)}: Znaczący podtyp algorytmu. Badacze zauważyli, że jeśli algorytm nie kwantyzuje macierzy losowo lub sekwencyjnie, ale w kolejności malejącej wariancji aktywacji (malejącej istotności macierzy Hesjana), to błędy zaokrągleń mogą zostać drastycznie zredukowane. Ten wariant określa się skrótowo jako \textbf{desc-act} (zazwyczaj w nazwie pliku modelu). Zwiększa on jednak czas kompresji dwu-trzykrotnie.
\end{itemize}

\section{AWQ (Activation-Aware Weight Quantization)}
\label{sec:awq}

Publicacja AWQ (Activation-Aware Weight Quantization) z 2023 roku wniosła do dziedziny nową perspektywę: perspektywę braku symetrii ważności poszczególnych kanałów modelu. Zauważono, że do tej pory metody (w tym GPTQ) analizowały błędy z perspektywy globalnej w przestrzeni wag modelu. AWQ dowodzi, że w architekturze modelu istnieje zaledwie około \textbf{0,1\%-1\% wag} które mają kluczowe znaczenie w procesie wnioskowania i modyfikowanie ich wywołuje katastrofalny spadek inteligencji modelu.

Te wagi nadrzędne są ściśle skorelowane ze zjawiskiem odpowiadającym im aktywacji o dużej amplitudzie. W przeciwieństwie do aktywacji odstających w BnB, w badaniach AWQ wykazano, że są to stałe cechy konkretnych kanałów, charakterystyczne dla typowych dystrybucji tekstowych; kanał, czy kolumna macierzy $W$ - otrzymuje zawsze silne pobudzenia w wektorze wejściowym $X$.

Podstawową różnicą AWQ względem GPTQ jest podejście do kalibracji. O ile ten pierwszy szuka optymalnego składnika błędu Hesjanowego, AWQ nie komplikuje sobie życia zaawansowaną algebrą liniową na poziomie szkolenia. Zamiast tego AWQ zauważa, że jeśli chcemy zapisać istotną wagę $w$ – to nawet duży błąd kwantyzacji $\Delta w$ przy niskiej intensywności aktywacji $x$ nie wywoła wielkich szkód w iloczynie wyjściowym. Jednakże – ten sam błąd wielkości $\Delta w$ przy bardzo dużej aktywacji $x$ spowoduje zniekształcenie wyjścia $x \cdot \Delta w$ całkowicie degradujące wynik. AWQ rozwiązuje to zjawisko stosując proste funkcjonalne \textbf{skalowanie (per-channel scaling)}. Tworzy wektor $s$ oraz szuka argumentu minimalizującego błąd replikacji $\mathrm{round}(w/s)$.

Dzięki takiemu przemyśleniu problemu, metoda AWQ potrafi drastycznie kompresować model analogicznie do GPTQ, jednak zachowując spójność dla tych 1% newralgicznych strumieni danych.

\subsection{Podtypy, warianty sprzętowe i architektoniczne}

Podobnie jak i GPTQ, AWQ rozwijane było jako projekt o otwartym kodzie źródłowym. Wyróżniamy w nim szereg wariantów wykonawczych i podtypów:

\begin{itemize}
    \item \textbf{Warianty bitowe W$X$A$Y$}: Architektury kwantyzowane AWQ są najczęściej opisywane notacją WxAy. Odnosi się to do precyzji wag oraz aktywacji. Bazowym i najpopularniejszym podtypem jest \textbf{W4A16} (4-bitowe wagi, 16-bitowe aktywacje). Autorzy zaproponowali również \textbf{W3A16} oraz eksperymentalne architektury \textbf{W8A8}. Ten wariant dotyczy problematycznej kwantyzacji samych wektorów wejściowych co jest obarczone bardzo dużym błędem maksymalnej wartości bezwzględnej.
    
    \item \textbf{Warianty kerneli (GEMM vs GEMV)}: Bardzo ważnym podziałem AWQ jest klasyfikacja ze względu na maszyny tensorowe na docelowym sprzęcie GPU. Są to tzw. Kernels w bibliotece AWQ:
    \begin{itemize}
        \item \textbf{GEMM-based (GEMM Kernel)}: Mnożenie macierzy zoptymalizowane pod pętle w architekturach Ampere (NVIDIA RTX 3000) i Hopper. AWQ zamiast dekwantyzować wagi na bieżąco do FP16, przekształca wagi do postaci rozpakowanej tuż przed wejściem do Tensor Core, opierając się o wzorce przekształceń GEMV. Format zalecany przy maksymalizacji liczby klatek na sekundę dla dużych rozmiarów partii (uruchamianie autorskich serwerów vLLM z włączonym stronnicowaniem).
        \item \textbf{GEMV-based (GEMV Kernel)}: General Matrix-Vector. Klasyczne dekwantyzacje przeprowadzane wyłącznie przy użyciu rejestrów CUDA. Stworzone specjalnie dla publikacji TinyChat i kart graficznych o bardzo małym standardzie obliczeniowym, takich jak mobilne układy Nvidia Jetson. Profile te są nastawione na wnioskowanie z rozmiarem partii, batch\_size=1. W ekosystemie jest to tzw. przypadek graniczny alokacji.
    \end{itemize}
    \item \textbf{Zoptymalizowane rozmiary grupy}: Podtypy AWQ zazwyczaj przyjmują wielkość bloku równą 128 (podobnie jak było to zaprezentowane w AutoGPTQ), co gwarantuje optymalny design na urządzeniach komórkowych. Wiele implementacji sprzętowych rezerwuje rozmiar grupy na 32 dla profilu wysokich architektur pro-użytkowych (np. procesory NPU).
\end{itemize}

\section{GGUF (GPT-Generated Unified Format)}
\label{sec:gguf}

Format GGUF (GPT-Generated Unified Format) to podstawowy fundament współczesnej inżynierii modeli na urządzenia brzegowe. Zaproponowany przez Georgiego Gerganowa (twórcę oprogramowania \texttt{llama.cpp}), wyparł on w szybkim czasie starszy format GGML. Motywacją do powstania GGUF był brak kompatybilności wstecznej i trudności w rozszerzaniu starych formatów o nowe architektury sieci, a także brak natywnego wsparcia dla zaawansowanych metod kwantyzacji.

GGUF definiuje strukturę przechowywania metadanych w eleganckiej strukturze typu klucz-wartość (KV) na początku pliku przechowującej parametry modelu, a tuż za nimi mapę wielkich tensorów reprezentującą wytrenowane, skwantyzowane parametry. Jest to format zoptymalizowany na poziomie C/C++ z natychmiastowym wsparciem operacji mapowania pamięci bez kopiowania (mmap), co ułatwia ładowanie modeli do pamięci RAM bez konwersji oraz wykorzystuje instrukcje SIMD (AVX2, AVX-512 na x86, NEON na architekturze ARM).

W przeciwieństwie do poprzednich metod (operujących na pojedynczych macierzach, jak AWQ), GGUF z założenia operuje na kompromisach sprzętowych dla poszczególnych warstw - to znaczy pozwala na ładowanie konkretnych warstw typu attention w różnej rozdzielczości w zależności od zapotrzebowania procesora.

\subsection{Architektura i mapa wartości GGUF}

Format GGUF jest niezwykle złożony, a wyselekcjonowane typy danych są indeksowane w sposób ciągły w celu łatwej interpretacji przez parser C++. Macierz pełnej mapy tensorów w formacie GGUF przedstawia się następująco:

\begin{table}[h]
\centering
\caption{Pełna mapa typów kwantyzacji w formacie GGUF}
\label{tab:gguf_map_full}
\begin{tabular}{|c|l|c|l|}
\hline
\textbf{Indeks} & \textbf{Format} & \textbf{Indeks} & \textbf{Format} \\ \hline
0 & F32 & 19 & IQ2\_XSS \\ \hline
1 & F16 & 20 & IQ2\_XS \\ \hline
2 & Q4\_0 & 21 & Q2\_K\_S \\ \hline
3 & Q4\_1 & 22 & IQ3\_XS \\ \hline
4 & Q4\_1\_SOME\_F16 & 23 & IQ3\_XXS \\ \hline
5 & Q4\_2 & 24 & IQ1\_S \\ \hline
6 & Q4\_3 & 25 & IQ4\_NL \\ \hline
7 & Q8\_0 & 26 & IQ3\_S \\ \hline
8 & Q5\_0 & 27 & IQ3\_M \\ \hline
9 & Q5\_1 & 28 & IQ2\_S \\ \hline
10 & Q2\_K & 29 & IQ2\_M \\ \hline
11 & Q3\_K\_S & 30 & IQ4\_XS \\ \hline
12 & Q3\_K\_M & 31 & IQ1\_M \\ \hline
13 & Q3\_K\_L & 32 & BF16 \\ \hline
14 & Q4\_K\_S & 33 & Q4\_0\_4\_4 \\ \hline
15 & Q4\_K\_M & 34 & Q4\_0\_4\_8 \\ \hline
16 & Q5\_K\_S & 35 & Q4\_0\_8\_8 \\ \hline
17 & Q5\_K\_M & 36 & TQ1\_0 \\ \hline
18 & Q6\_K & 37 & TQ2\_0 \\ \hline
 &  & 38 & MXFP4\_MOE \\ \hline
\end{tabular}
\end{table}

W tabeli tej można wyróżnić kilka rodzin rozwiązań, które często sprawiają trudności interpretacyjne. Powszechnie oznaczenia cyfrowe na pierwszej pozycji (\textit{Qx}) odnoszą się do docelowej pojemności kontenera (bity na wagę - bpw). Oznaczenia \textbf{"\_0" (np. Q4\_0, Q5\_0)} odnoszą się do modeli \textbf{z bezpośrednim skalowaniem}. Oznaczenia \textbf{"\_1" (np. Q4\_1, Q5\_1)} posiadają dodatkowo odchylenie minimum (tzw. przesunięcie (offset) dla każdego bloku, co daje większą rozdzielczość dynamiczną). Z kolei nowością sprzętową są podtypy indeksowe 33-35 (np. \textbf{Q4\_0\_4\_4}), które definiują modele określane jako "sub-block matrix" i dedykowane są one układom NPU w architekturze ARM gdzie instrukcje SIMD optymalizowane są na poziomie sprzętu w alokacjach pamięci L1/L2 4x4, 4x8 i 8x8 bajtów.

\begin{table}[h]
    \centering
    \caption{Wpływ kwantyzacji GGUF na rozmiar modelu na przykładzie Qwen3.6-27B}
    \label{tab:gguf_size_impact}
    \begin{tabular}{|c|c|}
        \hline
        \textbf{Format} & \textbf{Rozmiar Modelu (GB)} \\ \hline
        F16 & 55.6 \\ \hline
        IQ4\_XS & 15.4 \\ \hline
        IQ4\_NL & 16.1 \\ \hline
        Q3\_K\_S & 12.4 \\ \hline
        Q3\_K\_M & 13.6 \\ \hline
        Q4\_0 & 15.8 \\ \hline
        Q4\_1 & 17.3 \\ \hline
        Q4\_K\_S & 15.9 \\ \hline
        Q4\_K\_M & 16.8 \\ \hline
        Q5\_K\_S & 19.2 \\ \hline
        Q5\_K\_M & 19.5 \\ \hline
        Q6\_K & 22.5 \\ \hline
        Q8\_0 & 28.6 \\ \hline
    \end{tabular}
\end{table}


\subsection{Kwantyzacja typu K-Quants}
\textit{K-Quants} to zaimplementowana w GGUF metoda kwantyzacji podrzędnej, niejednorodnej na poziomie całego modelu. Bazuje na aproksymacji ważności tensorów w podziale na trzy poziomy wrażliwości: mały (Ma), średni (Mb) i duży (Mc). Rozwiązanie \textbf{K\_S (Small)} stosuje "najmocniejszą", najbardziej drastyczną formę redukcji. Generowane są tu profile Q2\_K, Q3\_K\_S.
Stosując podejście skwantyzowane do macierzy attention, okazało się, że tzw. mapa gradientowa jest niejednorodna, przez co powstały kolejne podrodziny:
\begin{itemize}
    \item \textbf{K\_M (Medium)}: Uznany zazwyczaj za kompromis rozwiązań progresywnych najmocniej kompresujących, a standardowym profilem. W przypadku podtypu Q4\_K\_M warstwy mniej istotnych tensorów oraz wyjść modelu są kompresowane w 4 bitach, ale newralgiczne bloki w modelach są kompresowane w wyższej rozdzielczości (6 bitów).
    \item \textbf{K\_L (Large)}: Profile z minimalnym uszczerbkiem na jakości. Utrzymują większość podmatryc na poziomie 6 lub 8 bitów, a kompresji do 4 bitów poddawane są jedynie ogony tensorów.
\end{itemize}

\subsection{Importance-Matrix oraz rodzina IQ}
Prawdziwym przełomem inżynierii wstecznej w metodzie GGUF jest \textbf{Importance Matrix} i podtypy \textbf{IQ} (Importance Quantization). W tej metodzie nie określamy istotności prosto z "wariancji aktywacji" lecz statystycznie z kalibracyjnego zbioru treningowego wyliczamy entropię punktową. Otrzymujemy n+k-wymiarowy tensor "ważności" dla wag, bazując na uogólnionym błędzie dowolnym.
Formaty takie jak \textbf{IQ1\_S, IQ1\_M} (indeksy 24, 31) umożliwiają agresywną kompresję do 1 bita na wagę (przez co wielkość pliku maleje ponad 16-krotnie!). Metody te są często znacznie lepsze od współczesnych podejść z rodziny klasycznego Q2 i Q3, właśnie ze względu na zgodność informacyjną z macierzą ważności.
Ostatnie aktualizacje GGUF z 2024 r. wprowadzają wsparcie dla podtypów architektur mixture-of-experts (MoE), takich jak \textbf{MXFP4\_MOE} (indeks 38) optymalizujących alokację routingów na poziomie bramek ekspertów.

\section{TurboQuant}
\label{sec:turboquant}

Na tle powyższych, wysoce rozwiniętych i sprofilowanych rozwiązań, TurboQuant stanowi nowatorskie rozwiązanie, które zdołało całkowicie namieszać w świecie optymizacji. Jest to metoda wywodząca się z dziedziny badań nad optymalizacją modeli i charakteryzuje się hybrydowym podejściem.

Problematyka operacji na liczbach całkowitych zakłada, że błędy skwantyzowanych prążków zwiększają się liniowo (zwiększa stratę sVRG) dla macierzy Hesjana wpisującego się w założenia rozszerzonych sieci rezydualnych, co przy modelach o skali 7B-70B powoduje tzw. cyfrową degradację wyjść. TurboQuant, zamiast korzystać ze skomplikowanej kalibracji z użyciem Hessiana jak GPTQ i AWQ, bezpośrednio zanurza wielowymiarową przestrzeń tensorów rzadkich do przestrzeni "gładkich" za sprawą transformacji Hadamarda i losowych macierzy ortogonalnych.

Główną ideą TurboQuant jest zaobserwowanie, że wstępnie skwantyzowana macierz wejścia według normy L2 w architekturze LLM ma strukturę quasi-izotropową. Zamanifestowało się to w filozofii losowych rotacji. Klastrowanie macierzy ortogonalnych i wykonanie tzw. szybkiej transformacji Walsha-Hadamarda (Fast Walsh-Hadamard Transform - FWHT) dla wejść przed włączeniem przybliżenia całkowitego, wytraca wagom ich wartości odstające w sposób prawie idealny. Z matematycznego punktu widzenia, TurboQuant nadaje macierzom "izotropię energetyczną" – dystrybucja jest spłaszczana. To usuwa anizotropię która prowadziła do dotychczasowych błędów. Fizycznie, usunięcie dużych oscylacji powoduje, że proste zaokrąglenie do INT4 czy nawet INT3 nie powoduje zapadania się klastrów odpowiadających za rozumowanie logiczne (Chain-of-Thought).

\subsection{Podtypy TurboQuant}

???

\begin{itemize}
    \item \textbf{TQ1\_0 oraz TQ2\_0}: Skwantyzowane z rzędu wektory istotnych cech zaimplementowanych do struktury GGUF (patrz Tabela \ref{tab:gguf_map_full}). W przeciwieństwie do klasycznych kwantyzacji 1-bitowych (jak binarne BNN), TQ2\_0 dzięki losowej rotacji gwarantuje, że w przestrzeni 2 bitów błąd kwantyzacji jest wysoce nieproporcjonalnie pomniejszony. Wdrożenia pokazały spadek jedynie o 1-2\% w testach MMLU i HumanEval w przeciwieństwie do katastrofalnych 20\% w starszych metodach 2-bitowych. Metoda TQ1\_0 jest absolutnym ekstremum kompresji, zaledwie 1 bit na wagę (wymaga dedykowanych kerneli dekompresujących).
    
    \item \textbf{TurboQuant-A (Approximate)}: Oparty na ucięciu (pruning) mało istotnych wektorów bazowych i transformacji Hadamarda w blokach 256x256. Umożliwia to wnioskowanie w czasie rzędu O(N log N).
    
    \item \textbf{Integracja z MoE}: Z racji wielkich osiągnięć TQ w utrzymaniu gładkiego rozkładu, metoda ta posłużyła jako fundament dla współczesnych MoE (Mixture of Experts) gdzie tradycyjny błąd kwantyzacji mógł zrujnować systemy routingu ekspertów. Podtyp \textbf{TQ-MoE} tworzy klastrową rotację dla ekspertów w siatce nawiązując bezpośrednio do formatów takich jak MXFP4\_MOE (indeks 38 w mapie GGUF). Pokazano tam, że TQ-MoE jest w stanie skompresować 8x7B model ekspertowy do reprezentacji 2-bitowej z lepszymi wynikami niż AWQ w 4-bit dla oryginalnych rozmiarów.
\end{itemize}

Ze względu na to, jak nowe rozwiązanie znacząco zmieniło kompresję (minimalizując wymóg kalibracji zbioru uczącego do niewielkiej liczby próbek, a często kalibrując w trakcie działania bez zbioru danych - zero-shot), TurboQuant stanowi dziś bardzo ważny punkt odniesienia w świecie vLLM i CUDA. Jego hybrydowe implementacje wymusiły na twórcach optymalizacje zunifikowane dla AVX-512.